%Version 3.1 December 2024
% See section 11 of the User Manual for version history
%
%%%%%%%%%%%%%%%%%%%%%%%%%%%%%%%%%%%%%%%%%%%%%%%%%%%%%%%%%%%%%%%%%%%%%%
%%                                                                 %%
%% Please do not use \input{...} to include other tex files.       %%
%% Submit your LaTeX manuscript as one .tex document.              %%
%%                                                                 %%
%% All additional figures and files should be attached             %%
%% separately and not embedded in the \TeX\ document itself.       %%
%%                                                                 %%
%%%%%%%%%%%%%%%%%%%%%%%%%%%%%%%%%%%%%%%%%%%%%%%%%%%%%%%%%%%%%%%%%%%%%

\documentclass[pdflatex,sn-mathphys-num]{sn-jnl}% Math and Physical Sciences Numbered Reference Style

%%%% Standard Packages
\usepackage{graphicx}%
\usepackage{multirow}%
\usepackage{amsmath,amssymb,amsfonts}%
\usepackage{amsthm}%
\usepackage{mathrsfs}%
\usepackage[title]{appendix}%
\usepackage{xcolor}%
\usepackage{textcomp}%
\usepackage{manyfoot}%
\usepackage{booktabs}%
\usepackage{algorithm}%
\usepackage{algorithmicx}%
\usepackage{algpseudocode}%
\usepackage{listings}%
%%%%

\theoremstyle{thmstyleone}%
\newtheorem{theorem}{Theorem}%
\newtheorem{proposition}[theorem]{Proposition}%

\theoremstyle{thmstyletwo}%
\newtheorem{example}{Example}%
\newtheorem{remark}{Remark}%

\theoremstyle{thmstylethree}%
\newtheorem{definition}{Definition}%

\raggedbottom

\begin{document}

\title[Multi-seed reproduction of NEU-ESC and ensemble distillation]{Reproducing NEU-ESC under a Multi-Seed Protocol and Improving It with Out-of-Fold Ensemble Distillation}

\author*{\fnm{Khang Lam} \sur{Huynh}\textsuperscript{*}}\email{khanglhse200666@fpt.edu.vn}

\author{\fnm{Dat Phan} \sur{Tuan}}\email{ptundatwork@gmail.com}

\author{\fnm{Huy Nguyen} \sur{Gia}}\email{huy522006@gmail.com}

\author{\fnm{Nhan Nguyen} \sur{Trong}}\email{hi.gesar@gmail.com}

\author{\fnm{Anh Vu} \sur{Ha}}\email{anhvh54@fpt.edu.vn}

\affil{\orgdiv{Faculty of Information Technology}, \orgname{FPT University}, \orgaddress{\city{Ho Chi Minh City}, \country{Vietnam}}}

\abstract{Student comments on university forums can show what learners need, and NEU-ESC is a recent Vietnamese dataset that labels such comments by sentiment and topic. Its authors report the best accuracy for multi-task BERT models with a SMART regularizer, but each setup has a single score and no spread, so small gains are hard to judge. We first reproduce the ten model rows from the public code, with 15 seeds per row. The mean absolute gap to the reported scores is 0.42 points, smaller than the mean seed range of 1.40 points, although no seed reaches four of the 20 reported cells. We then study Ensemble Distillation (ED). ED trains 30 fold models, averages their out-of-fold probabilities on the training set, and uses these averages as soft targets for one plain multi-task model. On NEU-ESC, ED beats all five reproduced baselines on both backbones and both tasks, in accuracy and macro-F1 (largest Holm-adjusted \textit{p} = 0.022). With PhoBERT it reaches 83.51 and 79.65, against 82.85 and 79.05 for MTL + SMART. Out-of-fold teachers give higher accuracy than in-sample teachers (+0.33 to +0.70), with no clear change in macro-F1. On UIT-VSFC, with the same settings and 10 seeds, ED improves topic accuracy over the strongest baseline, while its sentiment scores do not differ from that baseline after correction. ED keeps the inference cost of one model, but its teacher needs about 16 GPU hours on NEU-ESC, and its F1 on the rare Toxic class stays below the SMART setups.}

\keywords{Vietnamese text classification, Multi-task learning, Knowledge distillation, Reproducibility, Student feedback}

\maketitle

\section{Introduction}\label{sec:intro}

Comments on university forums show what students like and what they need, and recent work classifies such text in several languages \cite{malik2024student}, including Vietnamese \cite{phuc2026edupulse}. Vietnamese has strong pre-trained encoders such as PhoBERT \cite{nguyen2020phobert} and ViSoBERT \cite{nguyen2023visobert}, and later work extends this line to social media tasks such as hate speech detection \cite{nguyen2024vihatet5}. UIT-VSFC \cite{nguyen2018uitvsfc} holds over 16,000 student feedback sentences, and NEU-ESC \cite{mai2025neuesc} holds 32,966 forum comments with informal language, four sentiment classes and ten topic classes.

Mai et al. \cite{mai2025neuesc} treat the two labels of NEU-ESC as a multi-task learning (MTL) problem, in which one encoder feeds one output head per task \cite{caruana1997multitask,chen2024mtlnlp}. Their best sentiment score (83.66\%) comes from MTL with an auxiliary masked language modeling (MLM) head and SMART regularization \cite{jiang2020smart}, and their best topic score (79.80\%) comes from MTL with SMART. They weight the task losses equally, a simple choice that can match more complex multi-task optimizers \cite{royer2023scalarization}.

These results leave two problems open. The first is thin evidence. Each row of the original results table is a single score with no spread, and the public code sets no random seed. A single score is weak evidence, because the seed alone can move results by a statistically significant amount \cite{reimers2017reporting}, repeated runs differ on a fixed test set \cite{jordan2024variance}, and machine learning results are often hard to repeat \cite{semmelrock2025reproducibility}. Without repeated runs, a gain of half a point on NEU-ESC cannot be judged.

The second problem is prediction variance. In our 150 baseline models per task, 44.7\% of test comments are classified correctly by some models and wrongly by others, in each of the two tasks. Averaging the probabilities of 30 models reaches 84.14\% and 80.46\%, while one of those models reaches 82.09\% and 78.10\% on average. An ensemble of this size is too costly to serve, so we want one model that keeps part of this gain. Knowledge distillation trains one model on the outputs of a larger model or an ensemble \cite{bucilua2006compression,hinton2015distilling,moslemi2024survey}. Theory explains how the gain of an ensemble that differs only in random seeds can move into a single model \cite{allenzhu2023understanding}, while experiments show that students often do not match their teachers \cite{stanton2021distillation}. Recent work studies how logits should be matched \cite{sun2024logit}, and soft targets have also been studied in text classification \cite{ren2024smoothing}. We did not find a study that tests distillation from an out-of-fold ensemble on Vietnamese multi-task classification against baselines reproduced over many seeds.

This paper asks how close the public NEU-ESC code comes to the reported results over many runs, and whether a single plain MTL model can learn enough from an ensemble to beat every reproduced baseline. We make three contributions. First, we reproduce the ten BERT rows with 15 seeds each and test possible causes for the remaining gaps. Second, we evaluate Ensemble Distillation (ED), a plain MTL model trained with soft targets from 30 out-of-fold models. Distilling from an ensemble is not a new idea, and the new part is the careful test on this task, including an ablation of the out-of-fold step. Third, we check generalization on UIT-VSFC with the same settings, which also shows where ED does not help. The code, per-run results and report tables are available at \url{https://github.com/AI-Peak/neu-esc-reproduction}.

\section{Materials and methods}\label{sec:methods}

\subsection{Datasets}\label{subsec:data}

NEU-ESC was collected from university Facebook groups and online forums, and its authors cleaned the text and expanded common acronyms before release. UIT-VSFC contains student feedback sentences, which are shorter and fall into fewer classes. We use the published train, validation and test split of each dataset without change. Table~\ref{tab:data} gives the sizes and class shares. PhoBERT expects word-segmented input, while ViSoBERT reads raw text. For NEU-ESC, PhoBERT reads the segmented files published by the authors. For UIT-VSFC, we segment the text with underthesea 9.5.0, which differs from the VnCoreNLP tool commonly used with PhoBERT. All inputs are truncated to 128 tokens.

\begin{table}[!htb]
\caption{Datasets and class shares on the training split}\label{tab:data}
\footnotesize
\begin{tabular*}{\textwidth}{@{\extracolsep\fill}lp{0.85\textwidth}@{}}
\toprule
Task & Classes (share of training rows, \%) \\
\midrule
\multicolumn{2}{@{}l}{\textit{NEU-ESC (train 23,048, validation 3,305, test 6,613)}} \\
Sentiment & Neutral 69.1, Negative 15.7, Positive 12.6, Toxic 2.5 \\
Topic & Other 43.7, Academic 31.9, Service 7.1, Personal Affairs 4.5, News 2.7, Jobs \& Recruitment 2.4, Social Affairs 2.3, Help \& Share 2.0, Club \& Events 2.0, Spam 1.2 \\
\multicolumn{2}{@{}l}{\textit{UIT-VSFC (train 11,426, validation 1,583, test 3,166)}} \\
Sentiment & Positive 49.4, Negative 46.6, Neutral 4.0 \\
Topic & Lecturer 71.5, Training program 19.3, Others 4.9, Facility 4.3 \\
\botrule
\end{tabular*}
\end{table}

\subsection{Baseline models and reproduction}\label{subsec:baseline}

The baselines follow the public code of Mai et al. (commit b5637cc0), using PhoBERT-base-v2 and ViSoBERT at fixed revisions. The MTL model passes the [CLS] vector through a linear layer, a SiLU activation and layer normalization, and then sends it to one linear head per task. We reproduce five setups per backbone. Single-task trains one classifier per task for up to 15 epochs with early stopping. MTL uses the loss $0.5\,L_{\mathrm{sent}} + 0.5\,L_{\mathrm{topic}}$ for 20 epochs. MTL + SMART adds the SMART term with weight 0.02. MTL + MLM adds a masked token prediction head, and MTL + MLM + SMART uses both. All setups use AdamW with learning rate $2\times10^{-5}$, weight decay 0.01, batch size 128 and a linear learning rate decay.

We keep the authors' code where it is unusual, because the goal is to test what the public code produces. The most important case is SMART. The original method perturbs embeddings \cite{jiang2020smart}, but the public code adds noise to token IDs and casts the result back to integers. In practice this replaces some tokens with nearby IDs and asks the model to predict the same output, which is a form of consistency regularization. We keep this behavior and the name SMART. We also made a few recorded changes. We train with 16-bit floats on Kaggle T4 GPUs and split each batch into micro-batches. The MLM trainer always applies SMART, so we switch it off for the MTL + MLM row, and we limit random token IDs to the ViSoBERT vocabulary because the original range causes an error there. The authors state no checkpoint selection rule, so we select on the validation set. We therefore call our results a reproduction from the public code, not a recovery of the original experiment.

\subsection{Proposed method: Ensemble Distillation}\label{subsec:ed}

ED has two stages (Figure~\ref{fig:arch}). The first stage builds a fixed teacher, and the second stage trains one student with the same architecture as the MTL baseline. An ensemble lowers the variance between runs, so its outputs are a natural target, and a plain student keeps the inference cost of one model.

\begin{figure}[!htb]
\centering
\includegraphics[width=\textwidth]{figures/architecture.pdf}
\caption{Architecture of ED. Stage 1 builds soft targets from out-of-fold predictions, and Stage 2 trains the MTL student with cross-entropy and KL terms for both tasks. At inference time only the student runs. Output sizes are those of NEU-ESC}\label{fig:arch}
\end{figure}

In the first stage, we split the training set $\mathcal{D}$ into five fixed folds $\mathcal{D}_1,\dots,\mathcal{D}_5$, using a permutation seed that does not depend on the training seed. For each backbone $b$, training seed $s$ and fold $k$, we train an MTL model on $\mathcal{D}\setminus\mathcal{D}_k$, select its checkpoint on the validation set, and predict the rows of $\mathcal{D}_k$. Each training comment is therefore predicted only by models that did not see it, which is the same idea used in stacking \cite{wolpert1992stacked}. With two backbones, three seeds and five folds there are 30 models, and each comment receives six predictions. For a comment $x_i$ in fold $f(i)$ and a task $t$, the soft target is the mean probability
\begin{equation}
q_i^{t} = \frac{1}{|B|\,|S|}\sum_{b\in B}\sum_{s\in S} p^{t}_{b,s,f(i)}(x_i), \label{eq:teacher}
\end{equation}
where $B$ is the set of backbones and $S$ is the set of seeds. On NEU-ESC, these targets reach 83.55\% accuracy on the sentiment training labels and 80.47\% on the topic labels. The test set is not used at any point in this stage.

In the second stage, the student minimizes, from the first epoch,
\begin{equation}
L = \sum_{t\in\{\mathrm{sent},\,\mathrm{topic}\}} w_t \Big[ \mathrm{CE}\big(y^{t}, p^{t}\big) + \gamma\,\mathrm{KL}\big(q^{t} \,\|\, p^{t}\big) \Big], \label{eq:loss}
\end{equation}
with $w_t = 0.5$ and $\gamma = 1$. We use no temperature, so the probabilities are compared directly. If $H$ is entropy, each bracket in Eq.~\ref{eq:loss} can be rewritten as
\begin{equation}
\mathrm{CE}(y, p) + \gamma\,\mathrm{KL}(q \,\|\, p) = (1+\gamma)\,\mathrm{CE}\!\left(\frac{y+\gamma q}{1+\gamma},\, p\right) - \gamma H(q). \label{eq:mix}
\end{equation}
Since $H(q)$ does not depend on the student, ED with $\gamma = 1$ has the same gradients as cross-entropy against the average of the hard label and the teacher probabilities, scaled by two. Unlike uniform label smoothing, this soft target depends on the comment.

\section{Results}\label{sec:results}

\subsection{Experimental configurations and evaluation metrics}\label{subsec:config}

We use 15 seeds (42--56) per setup on NEU-ESC and 10 seeds (42--51) on UIT-VSFC. MTL checkpoints are selected by the mean validation macro-F1 of the two tasks, and single-task checkpoints by the validation macro-F1 of their task. Each selected model is evaluated once on the test set, and we report the mean and sample standard deviation over seeds, without choosing seeds or using test scores to pick an epoch. We report accuracy, macro-F1 (the plain mean of the per-class F1 scores) and weighted F1 (their mean weighted by class support). Comparisons use Welch's $t$-test. We apply Holm correction at $\alpha=0.05$ separately for each dataset, over every ED-versus-baseline pair in accuracy and macro-F1, which gives 40 tests on NEU-ESC and 32 on UIT-VSFC. Table~\ref{tab:insample} is corrected over its own eight tests. A difference is called significant when its adjusted $p$-value is below 0.05, unless we state otherwise.

\subsection{Reproduction of the NEU-ESC baselines}\label{subsec:repro}

Table~\ref{tab:repro} compares our 15-seed means with the reported scores. The mean absolute difference over the 20 cells is 0.42 points, while the range between the best and worst seed of a cell is 1.40 points on average (0.48 to 2.51). Our mean is below the reported score in 17 of 20 cells, as expected when a single score is compared with an average, and at least one seed reaches the reported score in 16 of 20 cells. Figure~\ref{fig:seeds} shows the spread of each setup and where each reported score falls in it.

\begin{table}[!htb]
\caption{Reproduction on NEU-ESC: test accuracy (\%, mean $\pm$ std over 15 seeds), reported score and difference ($\Delta$ = ours $-$ reported)}\label{tab:repro}
\footnotesize
\begin{tabular*}{\textwidth}{@{\extracolsep\fill}lcrrcrr@{}}
\toprule
 & \multicolumn{3}{c}{Sentiment} & \multicolumn{3}{c}{Topic} \\
\cmidrule{2-4}\cmidrule{5-7}
Setup & Ours & Reported & $\Delta$ & Ours & Reported & $\Delta$ \\
\midrule
\multicolumn{7}{@{}l}{\textit{PhoBERT}} \\
Single-task & $82.27\pm0.40$ & 81.75 & $+0.52$ & $78.16\pm0.58$ & 78.65 & $-0.49$ \\
MTL & $82.39\pm0.71$ & 82.88 & $-0.49$ & $78.75\pm0.36$ & 79.25 & $-0.50$ \\
MTL + SMART & $82.85\pm0.24$ & 83.26 & $-0.41$ & $79.05\pm0.17$ & 79.31 & $-0.26$ \\
MTL + MLM & $82.25\pm0.53$ & 82.67 & $-0.42$ & $78.46\pm0.44$ & 78.48 & $-0.02$ \\
MTL + MLM + SMART & $82.83\pm0.16$ & 83.66 & $-0.83$ & $79.09\pm0.17$ & 79.49 & $-0.40$ \\
\multicolumn{7}{@{}l}{\textit{ViSoBERT}} \\
Single-task & $82.30\pm0.64$ & 82.78 & $-0.48$ & $78.11\pm0.42$ & 77.94 & $+0.17$ \\
MTL & $82.04\pm0.43$ & 82.63 & $-0.59$ & $78.22\pm0.55$ & 78.56 & $-0.34$ \\
MTL + SMART & $82.71\pm0.18$ & 83.58 & $-0.87$ & $79.35\pm0.16$ & 79.80 & $-0.45$ \\
MTL + MLM & $81.83\pm0.68$ & 82.23 & $-0.40$ & $78.18\pm0.41$ & 78.38 & $-0.20$ \\
MTL + MLM + SMART & $82.82\pm0.16$ & 83.12 & $-0.30$ & $79.55\pm0.20$ & 79.25 & $+0.30$ \\
\botrule
\end{tabular*}
\end{table}

Four cells are never reached: ViSoBERT MTL + SMART (83.58 and 79.80) and PhoBERT MTL + MLM + SMART (83.66 and 79.49). The two sentiment cells are the largest gaps. Taking the epoch with the best test score and then the best seed reached only 83.09 and 83.25, so checkpoint choice does not explain them. Moving to 32-bit floats also did not help, with changes from $-0.28$ to $+0.41$ points on four setups. From the public material we could not find the cause of these two gaps.

The SMART effect depends on the metric. SMART beats MTL in accuracy in all four cells, as in the original paper. In macro-F1 it does not agree with the paper. SMART lowers macro-F1 in three of four cells, with a clear drop only for PhoBERT sentiment ($-0.79$, unadjusted $p=0.001$, where the paper reports $+1.53$), and it raises macro-F1 only for ViSoBERT topic ($+0.52$, against $+1.47$ in the paper).

\begin{figure}[!htb]
\centering
\includegraphics[width=\textwidth]{figures/fig_seed_distribution.pdf}
\caption{Test accuracy on NEU-ESC over 15 seeds per setup. Boxes show the quartiles and the median, whiskers reach the furthest seed within 1.5 times the interquartile range, and circles are seeds beyond that. Diamonds are the single scores reported by Mai et al. \cite{mai2025neuesc}, and the shaded column is ED}\label{fig:seeds}
\end{figure}

\subsection{Performance of ED on NEU-ESC}\label{subsec:ednue}

We first compare ED at the two highest points reported in the original paper. On PhoBERT sentiment, ED reaches $83.51\pm0.24$ against $82.83\pm0.16$ for the reproduced MTL + MLM + SMART (reported 83.66). On ViSoBERT topic, ED reaches $79.73\pm0.17$ against $79.35\pm0.16$ for the reproduced MTL + SMART (reported 79.80). These gains of 0.68 and 0.38 points have $p<0.001$ (unadjusted), and ED stays slightly below the two reported single scores, which no baseline seed reached.

Table~\ref{tab:all} gives every setup and metric. ED has the highest mean in all accuracy, macro-F1 and weighted F1 columns. All 40 Holm-corrected comparisons with the baselines are significant, and the largest adjusted $p$-value is 0.022 (ViSoBERT topic against MTL + MLM + SMART, in accuracy and in macro-F1). Against the strongest baseline of each column, the PhoBERT gains are $+0.66$ (sentiment accuracy), $+0.67$ (sentiment macro-F1), $+0.56$ (topic accuracy) and $+1.34$ (topic macro-F1), and the ViSoBERT gains are $+0.43$, $+0.71$, $+0.18$ and $+0.47$. In seven of these eight cells, even the worst ED seed is above the mean of the strongest baseline. ED also varies less between seeds than plain MTL, with an accuracy standard deviation of 0.17 to 0.27 against 0.36 to 0.71.

\begin{table}[!htb]
\caption{All setups on NEU-ESC (test, \%, mean $\pm$ std over 15 seeds; Weighted-F1 and F1 Toxic are means only). Bold marks the highest mean in each column and backbone}\label{tab:all}
\footnotesize
\setlength{\tabcolsep}{1.5pt}
\begin{tabular*}{\textwidth}{@{\extracolsep\fill}lccccccc@{}}
\toprule
 & \multicolumn{4}{c}{Sentiment} & \multicolumn{3}{c}{Topic} \\
\cmidrule{2-5}\cmidrule{6-8}
Setup & Accuracy & Macro-F1 & \begin{tabular}[c]{@{}c@{}}Weighted\\ F1\end{tabular} & \begin{tabular}[c]{@{}c@{}}F1\\ Toxic\end{tabular} & Accuracy & Macro-F1 & \begin{tabular}[c]{@{}c@{}}Weighted\\ F1\end{tabular} \\
\midrule
\multicolumn{8}{@{}l}{\textit{PhoBERT}} \\
Single-task & $82.27{\pm}0.40$ & $77.24{\pm}0.37$ & $82.27$ & 85.32 & $78.16{\pm}0.58$ & $61.49{\pm}0.58$ & $77.97$ \\
MTL & $82.39{\pm}0.71$ & $76.86{\pm}0.63$ & $82.39$ & 82.86 & $78.75{\pm}0.36$ & $62.23{\pm}0.41$ & $78.38$ \\
MTL + SMART & $82.85{\pm}0.24$ & $76.07{\pm}0.57$ & $82.02$ & \textbf{85.35} & $79.05{\pm}0.17$ & $61.95{\pm}0.60$ & $78.23$ \\
MTL + MLM & $82.25{\pm}0.53$ & $76.85{\pm}0.41$ & $82.32$ & 82.79 & $78.46{\pm}0.44$ & $61.92{\pm}0.54$ & $78.18$ \\
MTL + MLM + SMART & $82.83{\pm}0.16$ & $76.23{\pm}0.54$ & $82.08$ & \textbf{85.35} & $79.09{\pm}0.17$ & $62.11{\pm}0.60$ & $78.24$ \\
ED & $\mathbf{83.51{\pm}0.24}$ & $\mathbf{77.91{\pm}0.28}$ & $\mathbf{83.42}$ & 83.28 & $\mathbf{79.65{\pm}0.27}$ & $\mathbf{63.57{\pm}0.40}$ & $\mathbf{79.27}$ \\
\multicolumn{8}{@{}l}{\textit{ViSoBERT}} \\
Single-task & $82.30{\pm}0.64$ & $76.51{\pm}0.46$ & $82.18$ & 83.39 & $78.11{\pm}0.42$ & $61.32{\pm}0.41$ & $77.79$ \\
MTL & $82.04{\pm}0.43$ & $76.08{\pm}0.76$ & $81.93$ & 82.60 & $78.22{\pm}0.55$ & $61.70{\pm}0.60$ & $77.95$ \\
MTL + SMART & $82.71{\pm}0.18$ & $75.73{\pm}0.43$ & $81.88$ & 84.49 & $79.35{\pm}0.16$ & $62.21{\pm}0.41$ & $78.49$ \\
MTL + MLM & $81.83{\pm}0.68$ & $76.00{\pm}0.47$ & $81.79$ & 82.53 & $78.18{\pm}0.41$ & $61.56{\pm}0.47$ & $77.89$ \\
MTL + MLM + SMART & $82.82{\pm}0.16$ & $76.05{\pm}0.42$ & $82.07$ & \textbf{84.68} & $79.55{\pm}0.20$ & $62.58{\pm}0.67$ & $78.76$ \\
ED & $\mathbf{83.25{\pm}0.18}$ & $\mathbf{77.22{\pm}0.28}$ & $\mathbf{83.05}$ & 82.94 & $\mathbf{79.73{\pm}0.17}$ & $\mathbf{63.05{\pm}0.26}$ & $\mathbf{79.21}$ \\
\botrule
\end{tabular*}
\end{table}

Figure~\ref{fig:class} shows where the gain comes from, using MTL + SMART as the reference. Positive and Negative gain 2.7 to 4.7 points of F1, while Neutral, the largest class, changes by about 0.2. Toxic moves the other way ($-2.08$ and $-1.55$). Rare topics differ by backbone, with wide intervals.

\begin{figure}[!htb]
\centering
\includegraphics[width=\textwidth]{figures/fig_perclass_delta.pdf}
\caption{Change in per-class F1 on NEU-ESC, ED minus MTL + SMART. Points are mean differences over 15 seeds per setup, and lines are uncorrected 95\% Welch confidence intervals. Parentheses give the training share}\label{fig:class}
\end{figure}

\subsection{Ablation of the teacher}\label{subsec:ablation}

Table~\ref{tab:ablation} removes the KL term and compares ED with its own teacher. The KL term alone adds $+0.90$ to $+1.51$ points over plain MTL ($p<0.001$), more than SMART provides in each cell. A single ED model is above the 15-model ensemble of its own backbone in all four cells, although the members of that ensemble see only 80\% of the training data. ED closes 34\% to 51\% of the gap between MTL + SMART and the 30-model ensemble, at 1/30 of the inference cost.

\begin{table}[!htb]
\caption{Ablation on NEU-ESC: test accuracy (\%) without the KL term, and against the teacher. The gap closed is (ED $-$ MTL + SMART) / (30-model ensemble $-$ MTL + SMART)}\label{tab:ablation}
\footnotesize
\begin{tabular*}{\textwidth}{@{\extracolsep\fill}lcccc@{}}
\toprule
 & \multicolumn{2}{c}{PhoBERT} & \multicolumn{2}{c}{ViSoBERT} \\
\cmidrule{2-3}\cmidrule{4-5}
Setup & Sentiment & Topic & Sentiment & Topic \\
\midrule
MTL (ED without the KL term) & 82.39 & 78.75 & 82.04 & 78.22 \\
MTL + SMART & 82.85 & 79.05 & 82.71 & 79.35 \\
ED & \textbf{83.51} & \textbf{79.65} & \textbf{83.25} & \textbf{79.73} \\
Ensemble of 15 fold models, same backbone & 83.28 & 79.46 & 82.81 & 79.34 \\
Ensemble of all 30 fold models & 84.14 & 80.46 & 84.14 & 80.46 \\
Gap closed & 51\% & 43\% & 38\% & 34\% \\
\botrule
\end{tabular*}
\end{table}

To test the out-of-fold step, we also built an in-sample teacher from six MTL models (two backbones, three seeds) trained on the full training set, which then predict the same comments they learned. This teacher reaches 97.83\% and 96.03\% on the training labels, against 83.55\% and 80.47\% for the out-of-fold teacher, so its targets are close to the hard labels. In Table~\ref{tab:insample}, out-of-fold targets give significantly higher accuracy in all four cells ($+0.33$ to $+0.70$), but the macro-F1 differences are not significant. In-sample ED already gives 38\% to 73\% of the accuracy gain and 69\% to 101\% of the macro-F1 gain of out-of-fold ED over MTL. The two teachers also differ in data, because the in-sample models train on all of it and the fold models on 80\%.

\begin{table}[!htb]
\caption{Out-of-fold (OOF) against in-sample teachers on NEU-ESC (test, \%, 15 seeds). Holm correction is applied over these eight tests}\label{tab:insample}
\footnotesize
\setlength{\tabcolsep}{3pt}
\begin{tabular*}{\textwidth}{@{\extracolsep\fill}llcccccc@{}}
\toprule
 & & \multicolumn{3}{c}{Accuracy} & \multicolumn{3}{c}{Macro-F1} \\
\cmidrule{3-5}\cmidrule{6-8}
Backbone & Task & In-sample & OOF & Holm $p$ & In-sample & OOF & Holm $p$ \\
\midrule
PhoBERT & Sentiment & $82.82\pm0.37$ & $83.51\pm0.24$ & $<0.001$ & $77.58\pm0.40$ & $77.91\pm0.28$ & 0.061 \\
PhoBERT & Topic & $79.09\pm0.28$ & $79.65\pm0.27$ & $<0.001$ & $63.16\pm0.48$ & $63.57\pm0.40$ & 0.061 \\
ViSoBERT & Sentiment & $82.93\pm0.21$ & $83.25\pm0.18$ & $<0.001$ & $77.11\pm0.40$ & $77.22\pm0.28$ & 0.726 \\
ViSoBERT & Topic & $79.25\pm0.22$ & $79.73\pm0.17$ & $<0.001$ & $63.07\pm0.42$ & $63.05\pm0.26$ & 0.876 \\
\botrule
\end{tabular*}
\end{table}

\subsection{Generalization to UIT-VSFC}\label{subsec:uit}

We repeat the full pipeline on UIT-VSFC with the same hyperparameters and no tuning, and we rebuild the teacher from its training set. Single-task was not run here, so there are four baselines (Table~\ref{tab:uit}). Of the 32 Holm-corrected comparisons, 21 are significant, and ED is significantly below a baseline in none. Against the strongest baseline of each column, topic accuracy rises by 0.60 points on PhoBERT and 0.41 on ViSoBERT (adjusted $p\le0.002$), and topic macro-F1 rises by 0.90 on ViSoBERT (adjusted $p=0.011$). The PhoBERT macro-F1 gain of 0.69 does not survive correction (adjusted $p=0.131$). Plain MTL has a wide seed spread here, so the topic accuracy gains over MTL also fall short after correction (adjusted $p=0.127$ and 0.065).

\begin{table}[!htb]
\caption{Results on UIT-VSFC (test, \%, mean $\pm$ std over 10 seeds). Bold marks the highest mean in each column and backbone}\label{tab:uit}
\footnotesize
\begin{tabular*}{\textwidth}{@{\extracolsep\fill}lccccc@{}}
\toprule
 & \multicolumn{3}{c}{Sentiment} & \multicolumn{2}{c}{Topic} \\
\cmidrule{2-4}\cmidrule{5-6}
Setup & Accuracy & Macro-F1 & F1 Neutral & Accuracy & Macro-F1 \\
\midrule
\multicolumn{6}{@{}l}{\textit{PhoBERT}} \\
MTL & $93.77\pm0.16$ & $82.11\pm0.82$ & 55.3 & $89.05\pm0.70$ & $80.01\pm0.62$ \\
MTL + SMART & $92.44\pm0.20$ & $76.86\pm0.99$ & 41.9 & $88.96\pm0.22$ & $78.78\pm0.72$ \\
MTL + MLM & $93.84\pm0.13$ & $\mathbf{82.59\pm0.50}$ & \textbf{56.6} & $89.10\pm0.40$ & $79.86\pm0.71$ \\
MTL + MLM + SMART & $92.54\pm0.21$ & $77.14\pm0.95$ & 42.5 & $89.11\pm0.23$ & $79.13\pm0.46$ \\
ED & $\mathbf{93.88\pm0.12}$ & $82.02\pm0.51$ & 54.9 & $\mathbf{89.71\pm0.11}$ & $\mathbf{80.70\pm0.58}$ \\
\multicolumn{6}{@{}l}{\textit{ViSoBERT}} \\
MTL & $92.71\pm0.20$ & $80.47\pm1.17$ & 52.2 & $88.77\pm0.48$ & $79.47\pm0.43$ \\
MTL + SMART & $92.36\pm0.24$ & $78.11\pm0.64$ & 46.0 & $88.91\pm0.19$ & $79.52\pm0.58$ \\
MTL + MLM & $92.80\pm0.29$ & $\mathbf{80.71\pm0.87}$ & \textbf{52.8} & $88.54\pm0.48$ & $79.27\pm0.69$ \\
MTL + MLM + SMART & $92.42\pm0.24$ & $78.32\pm0.66$ & 46.5 & $88.85\pm0.25$ & $79.43\pm0.65$ \\
ED & $\mathbf{93.08\pm0.10}$ & $80.34\pm0.49$ & 51.3 & $\mathbf{89.32\pm0.16}$ & $\mathbf{80.42\pm0.34}$ \\
\botrule
\end{tabular*}
\end{table}

On sentiment, ED has the highest accuracy, but it does not differ significantly from MTL + MLM (PhoBERT $+0.03$, adjusted $p=1$; ViSoBERT $+0.28$, adjusted $p=0.121$). MTL + MLM has the best sentiment macro-F1, and ED is lower by 0.57 and 0.37 points, also not significant after correction (adjusted $p=0.135$ and 1). Against MTL + SMART, ED wins all eight accuracy and macro-F1 comparisons.

Across the eight accuracy cells of both datasets, the gain of ED over plain MTL grows with the gap between the teacher and plain MTL (Pearson $r=0.94$, $n=8$). The cells share teachers and are not independent, so this number is a description and not a test. In the clearest case, PhoBERT sentiment on UIT-VSFC, plain MTL (93.77) already matches the 30-model ensemble (93.75), and ED gains only 0.11 points.

\subsection{Rare classes, logit bias and cost}\label{subsec:rare}

ED trades some rare-class F1 for its overall gain. On NEU-ESC, its Toxic F1 is below single-task and the SMART setups, and above MTL and MTL + MLM (Table~\ref{tab:all}). On UIT-VSFC, its Neutral F1 is level with MTL ($p>0.4$), below MTL + MLM, and above the SMART setups (Table~\ref{tab:uit}). The teacher itself under-predicts these classes. Its most likely class is Toxic for 2.04\% of NEU-ESC training rows, while 2.54\% carry that label, and it picks Neutral for 2.92\% of UIT-VSFC rows against a label share of 4.01\%.

A simple way to shift a decision toward a rare class is post-hoc logit adjustment \cite{menon2021logit}. We add one bias to the logit of the rare sentiment class, choose it on validation macro-F1 from a fixed grid, and apply it once to the test set. On UIT-VSFC this raises Neutral F1 by 2.26 and 4.80 points and macro-F1 by 0.75 and 1.63 points (paired $p\le0.019$), with no significant change in accuracy, and ED is then no longer below MTL + MLM in macro-F1 (PhoBERT $+0.03$, $p=0.905$; ViSoBERT $+0.59$, $p=0.052$). On NEU-ESC, Toxic F1 does not change significantly ($-0.07$, $p=0.754$; $+0.33$, $p=0.161$) and it stays below the best SMART setup under the same treatment (83.21 against 85.40 on PhoBERT, 83.26 against 85.37 on ViSoBERT). These $p$-values are not corrected, so the main results of this paper stay without bias.

\textit{Cost.} A single ED run takes 0.65 hours on NEU-ESC, the same as plain MTL and about 2.7 times faster than MTL + SMART (1.72 hours). The extra cost is the teacher, once per dataset: 15.65 GPU hours on NEU-ESC (3.89 for the in-sample teacher) and 3.79 on UIT-VSFC. Counting the teacher, the first ED model on NEU-ESC needs 16.30 hours, against 0.65 to 2.19 hours for the MTL baselines.

\section{Discussion}\label{sec:discussion}

The public NEU-ESC code reproduces most of the reported scores within seed noise, and ED beats all reproduced baselines on NEU-ESC, with a smaller effect on UIT-VSFC.

A likely reason for the gain is the combination of unstable predictions and informative soft targets. Nearly half of the test comments in NEU-ESC change between correct and wrong across models, and Eq.~\ref{eq:mix} shows that ED trains on a per-comment mixture of the label and the ensemble average, which may reduce the part of the error that depends on the seed. The lower seed spread of ED agrees with this idea (unadjusted Brown--Forsythe $p<0.05$ against MTL in three of four NEU-ESC accuracy cells), but we did not test the explanation directly. Our results agree with theory on moving ensemble gains into one model \cite{allenzhu2023understanding}, and with the observation that a student often does not match its teacher \cite{stanton2021distillation}, since ED stays below the 30-model ensemble in all NEU-ESC cells.

The ablation adds a caution to the common argument for held-out targets. Out-of-fold targets raise accuracy, yet most of the macro-F1 gain is already there with in-sample targets, so the choice of teacher depends on the goal.

SMART, as implemented in the public code, raises accuracy on NEU-ESC but lowers macro-F1 in most cells, and on UIT-VSFC it loses in nearly every comparison. A single reported score would not reveal this pattern \cite{reimers2017reporting}. For practice, ED fits a case that needs a higher overall score from one served model. When Toxic F1 is the main goal on NEU-ESC, single-task or a SMART setup is still the better choice, and on UIT-VSFC a validated logit bias can recover Neutral F1.

Our study has limits. The $p$-values describe training randomness on one fixed test set and do not cover the sampling of test comments, and run-to-run differences on a test set can overstate differences at the distribution level \cite{jordan2024variance}, so small gains such as $+0.18$ need care. We fixed $\gamma=1$ and the teacher size in advance without tuning. Our reproduction is not an exact copy of the original experiment, because of the recorded changes in Section~\ref{subsec:baseline} and the two reported cells that remain out of reach. We also suspect that the student inherits the rare-class under-prediction of its teacher, but we did not test it.

\section{Conclusion}\label{sec:conclusion}

We measured how well the NEU-ESC results hold over many seeds, and whether one plain multi-task model can learn enough from an ensemble to beat them. Over 15 seeds per row, the public code comes within 0.42 points of the reported scores on average, and two sentiment scores (83.58 and 83.66) remain out of reach even with test-based checkpoint choice. ED beats every reproduced baseline on NEU-ESC in accuracy and macro-F1 after Holm correction, with the inference cost of plain MTL. On UIT-VSFC, ED improves topic accuracy over the strongest baseline, and sentiment does not differ from it. Out-of-fold targets add accuracy over in-sample targets, while most of the macro-F1 gain comes from distillation itself.

ED has two clear costs: about 16 GPU hours for the teacher on NEU-ESC, and a lower Toxic F1 than the SMART setups. Our next experiments are a lower $\gamma$ for rare classes, class-balanced losses with ED, $\gamma \in \{0.5, 2\}$, and teachers of 5, 15 and 30 models.

\backmatter

\section*{Declaration}

Declaration of generative AI and AI-assisted technologies in the writing process: Claude was used for assisting with language editing only. The authors reviewed and revised the manuscript and take full responsibility for the accuracy and integrity of the work.

\bibliography{references}

\end{document}
